\section{Evidence Index for the Assessment Criteria}
\label{app:evidence}

This appendix maps the report's principal claims to their supporting
artefacts and reproduction commands. It is an index of evidence, not an
additional experiment or a declaration that every submission item is complete.
All file paths below are relative to the repository's \path{project2/}
directory; commands in D3 also start there.

The unified \path{code/run_all.py} regenerates all nine external report
figures from archived measurements and six neural checkpoints. With
\texttt{--recompute}, it also reruns the classical experiments and writes
fresh CSV files and four tables under \path{algorithm_section/}. The report's
checked-in tables retain the measurements stated in their captions; new
hardware timings must be reviewed before replacing those snapshots.
The architecture figure is native TikZ, rebuilt with the report.
The interface-order and stability scripts in \path{report/support/verify/}
use separate small models; the matrix-evidence script deliberately imports
the production assembler and provides numerical corroboration.

\subsection*{D1. Criteria mapped to evidence}

The letters below are the criteria of the course assessment scheme; the column
that matters for a reader is the third one, which names the artefact in this
report that carries each claim.

{\footnotesize
\begin{longtable}{@{}>{\raggedright\arraybackslash}p{2.55cm}>{\raggedright\arraybackslash}p{3.85cm}>{\raggedright\arraybackslash}p{4.75cm}>{\raggedright\arraybackslash}p{3.55cm}@{}}
\caption{Mapping of the assessment criteria to the evidence in this report. Column 2 names
the claim, column 3 the artefact that supports it, column 4 how to regenerate it.}\label{tab:evidence_index}\\
\toprule
Criterion & Claim it assesses & Evidence in this report & Regeneration \\
\midrule
\endfirsthead
\caption[]{Mapping of the assessment criteria to evidence (continued).}\\
\toprule
Criterion & Claim it assesses & Evidence in this report & Regeneration \\
\midrule
\endhead
\bottomrule
\endfoot
A. Mathematical soundness
  & The formulation, the conservative stencil and the discrete structure are
    derived, not asserted.
  & \S1.2 model, weak form and Lax--Milgram well-posedness (Theorem~\ref{thm:well_posed});
    \S2.1 the flux balance (eq.~\ref{eq:discrete_balance}) leading to the
    stencil (eq.~\ref{eq:conservative_fdm}); \S2.2 second-order truncation error
    (Theorem~\ref{thm:truncation_smooth}); \S2.4 SPD/$M$-matrix/DMP
    (Theorem~\ref{thm:spd_mmatrix}) and the stability--convergence theorem
    (Theorem~\ref{thm:convergence}).
  & analytic; no code needed \\[2pt]
A. Mathematical soundness
  & The derivation of symmetry and positive definiteness has numerical corroboration.
  & Table~\ref{tab:matrix_evidence}: $\|A-A^\top\|/\|A\|=0$ in both norms; the
    smallest-eigenvalue estimates for five cases lie in $[28.675,48.849]$, above $8a_{\min}=8$; dense
    Cholesky succeeds in all five cases.
  & \texttt{check\_matrix\_}\allowbreak\texttt{evidence} (see D3) \\[2pt]
A. Mathematical soundness
  & The claimed order is observed, in two named norms, on four meshes.
  & Figure~\ref{fig:alg_accuracy}: fits on the finest three meshes give orders $1.99984$ (FDM) and $1.99939$
    (FEM) in $L^\infty$, $1.99996$ and $1.99898$ in the common integrated
    $L^2$.
  & \texttt{code/run\_all.py} $\to$ \path{algorithm_section/results/} \\[2pt]
B. Depth and novelty
  & The two discretisations are compared under matched settings, with a cost
    measure.
  & \S4.2 and Table~\ref{tab:alg_smooth} (equal unknowns $=(N-1)^2$, same
    solver tolerance); Table~\ref{tab:alg_cost}: assembly
    $2.738\times10^{-3}$~s (FDM) versus $1.026\times10^{-1}$~s (FEM) at
    $N=128$, with $n_{\mathrm{nz}}=80137$ in both.
  & \texttt{code/run\_all.py} \\[2pt]
B. Depth and novelty
  & An additional extension beyond the prescribed comparison exists and is
    reported honestly.
  & \S\ref{sec:interface_pinn}: interface-decomposed PINN against harmonic FDM
    and fitted FEM, Table~\ref{tab:pinn_accuracy} ($E_{\mathrm{rel,samp}}=3.12\times10^{-5}$,
    maximum interface-flux error $1.37\times10^{-3}$), the single-network ablation in Table~\ref{tab:pinn_ablation}, and the
    explicit distinction from the published I-PINNs architecture.
  & the PINN replay script (D3) \\[2pt]
C. Implementation
  & The submitted code produces the submitted numbers, and can be run.
  & \S\ref{sec:algorithm_impl} documents the data layout and the assembly;
    Figure~\ref{fig:alg_spy} shows the actual matrix pattern against the
    hand-check matrix; the residual histories (Figure~\ref{fig:alg_solver}) and tolerance sweep
    (Table~\ref{tab:alg_tol}) are produced by the
    same functions the report's \S2 describes.
  & \texttt{python -B code/run\_all.py} \\[2pt]
D. Experiments and visualisation
  & The experiments answer the question asked, not a nearby one.
  & \S4.3 and Table~\ref{tab:interface_refinement}: the layered benchmark is
    reported with the exact solution, the shared-face flux, the two-sided
    reconstructed jump ($J_q$) and the error against the exact flux $J_*$,
    instead of a single jump statistic; Figure~\ref{fig:alg_interface} plots
    the solution and the numerical normal flux across $x=1/2$.
  & \texttt{code/run\_all.py} \\[2pt]
D. Experiments and visualisation
  & The claimed maximum-principle property is tested, not just proved.
  & \S4.3 reports that both arithmetic and harmonic solutions stay in $[0,1]$,
    which is what the sign structure of Theorem~\ref{thm:spd_mmatrix} predicts;
    the bound is stated as a matrix property and is not used as accuracy
    evidence.
  & \texttt{code/run\_all.py} $\to$ \path{algorithm_section/results/} \\[2pt]
D. Experiments and visualisation
  & Conditioning and iterative cost are measured under refinement \emph{and}
    contrast, at comparable residual tolerances.
  & Table~\ref{tab:conditioning_corrected}: $\kappa_2$ rises
    $1.61\times10^{3}\to7.15\times10^{3}\to6.49\times10^{4}$ at $N=63$ for
    $c=1,10,100$; CG iterations $216\to521\to1354$; Jacobi-CG $216/224/226$.
    All use requested CG rtol $10^{-12}$, with actual residuals recorded.
  & \texttt{code/run\_all.py} $\to$ \path{algorithm_section/results/conditioning.csv} \\[2pt]
E. Writing and AI transparency
  & The algebraic error is separated from the discretisation error, and the
    tolerance used is stated.
  & \S4.1 and Table~\ref{tab:alg_tol}: at $N=32$ the CG solution agrees with
    sparse LU to $5.33\times10^{-15}$ (FDM) and $1.01\times10^{-14}$ (FEM) at
    $\mathrm{rtol}=10^{-12}$, far below the discretisation error
    $8.0\times10^{-4}$; $\delta_{\mathrm{alg}}/E_\infty<10^{-3}$ on all four
    meshes.
  & \texttt{code/run\_all.py} \\[2pt]
E. Writing and AI transparency
  & AI assistance is disclosed, audited and attributable.
  & Appendix~\ref{app:ai_log}: tool table, claim audits with the falsifier
    that decided each one, the independent check and its evidence location, and
    the ownership declaration; the audit extract and its SHA-256 live in
    \path{report/support/ai_log/}. Separate contributor records in Appendix~\ref{app:ai_log} document
    Tianrui Ji's algorithm contribution and Han Liu's writing contribution.
  & \path{code/check_artifacts.py} (D3) \\[2pt]
F. Individual contribution
  & Contributions are attributed to individual members and linked to artefacts.
  & Appendix~\ref{app:ics} contains five labelled statements. The contributor
    logs identify code, sections and trusted library interfaces; the ownership
    table in \path{report/APPENDIX/AI_Theory_MeichenDing.tex} contains Meichen
    Ding's row only. Signatures and draft confirmations remain individual tasks.
  & \texttt{git log --author} (repository) \\
\end{longtable}
}

\subsection*{D2. Headline numbers and where they come from}

{\footnotesize
\begin{longtable}{@{}>{\raggedright\arraybackslash}p{5.55cm}>{\raggedright\arraybackslash}p{3.05cm}>{\raggedright\arraybackslash}p{2.35cm}>{\raggedright\arraybackslash}p{3.75cm}@{}}
\caption{The numbers quoted in the abstract, \S\ref{sec:results} and
\S\ref{sec:interface_pinn}, with the artefact that produces each one.}\label{tab:evidence_numbers}\\
\toprule
Quantity & Value & Stated in & Produced by \\
\midrule
\endfirsthead
\caption[]{Headline numbers and their provenance (continued).}\\
\toprule
Quantity & Value & Stated in & Produced by \\
\midrule
\endhead
\bottomrule
\endfoot
Symmetry defect $\|A-A^\top\|/\|A\|$ & $0$ (max and Frobenius) & Table~\ref{tab:matrix_evidence} & \path{report/support/verify/} \\
Smallest eigenvalue $\lambda_{\min}$, smooth & $28.675$ / $28.694$ ($N=32,64$) & Table~\ref{tab:matrix_evidence} & \texttt{eigsh}, \texttt{which='SA'} \\
Fitted order, $L^\infty$, smooth & $1.99984$ (FDM), $1.99939$ (FEM) & \S4.2 & \path{algorithm_section/results/smooth.csv} \\
Forcing defect of the manufactured source & $7.11\times10^{-15}$ & \S4.1 & complex-step difference \\
Affine-patch check of the Dirichlet elimination & $4.89\times10^{-15}$ & \S4.1 & dense reference solve \\
Harmonic FDM, layered $c=10$, $N=63$ & $E_\infty=7.33\times10^{-15}$, $E_q=1.08\times10^{-13}$ & Table~\ref{tab:interface_refinement} & \path{test/results/interface_layered.csv} \\
Arithmetic FDM, layered $c=10$, $N=63$ & $E_\infty=9.60\times10^{-3}$, $E_q=0.662$, $J_q=3.85\times10^{-13}$ & Table~\ref{tab:interface_refinement} & \path{algorithm_section/results/layered.csv} \\
Arithmetic nodal orders, $N=31\to255$ & $0.991$, $0.995$, $0.998$ & \S4.3 & \path{algorithm_section/results/layered.csv} \\
Midline interface-flux error, arithmetic & $4.70\times10^{-2}\to5.64\times10^{-3}$ & \S4.3 & \path{test/results/interface_layered.csv} \\
Entrywise ratio of the two inverses & $3.1024>1$ (counterexample holds) & Remark~\ref{rem:entrywise_false} & \path{report/support/verify/check_stability.py} \\
$\|A^{-1}\|_\infty$ against $1/(8a_{\min})$ & $0.0861$ (1D), $0.0502$ (2D); bound $0.125$ & Theorem~\ref{thm:convergence}, Remark~\ref{rem:conv_constant} & \path{report/support/verify/check_stability.py} \\
CG\,--\,LU discrepancy at $\mathrm{rtol}=10^{-12}$ & $5.33\times10^{-15}$ (FDM), $1.01\times10^{-14}$ (FEM) & Table~\ref{tab:alg_tol} & \path{report/tables/tolerance.tex} \\
Cost at $N=128$ (assembly / CG / LU) & $2.74/15.3/30.6$~ms (FDM), $103/34.7/30.8$~ms (FEM) & Table~\ref{tab:alg_cost} & \path{report/tables/cost.tex} \\
Condition number at $c=1,10,100$, $N=63$ & $1.61\times10^{3}$, $7.15\times10^{3}$, $6.49\times10^{4}$ & Table~\ref{tab:conditioning_corrected} & \path{algorithm_section/results/conditioning.csv} \\
Interface PINN parameters & $17{,}154$ & Table~\ref{tab:pinn_setup} & archived config \\
Interface PINN accuracy, $c=10$ & $E_{\mathrm{rel,samp}}=3.12\times10^{-5}$; flux error $1.37\times10^{-3}$ & Table~\ref{tab:pinn_accuracy} & \path{report/support/interface_pinn/} \\
Single-network ablation & $E_{\mathrm{rel,samp}}=0.310$ versus $3.12\times10^{-5}$; unequal parameter counts & Table~\ref{tab:pinn_ablation} & \path{report/support/interface_pinn/comparison.json} \\
Interface PINN training time & $231.3\pm4.9$~s (three seeds) & Table~\ref{tab:pinn_endpoints} & archived logs \\
\end{longtable}
}

\subsection*{D3. Reproduction commands}

\begin{verbatim}
# Start in project2; replay archived evidence and regenerate every figure:
python -B code/run_all.py
# Optional: re-solve classical experiments (replaces their timings):
python -B code/run_all.py --recompute
# Optional: also build report/Method.pdf; requires pdfLaTeX/BibTeX:
python -B code/run_all.py --build-report

# Targeted theory checks:
python -B report/support/verify/check_interface_order.py
python -B report/support/verify/check_stability.py
python -B report/support/verify/check_matrix_evidence.py
# Check vector PDFs, 300-dpi previews, archive and AI-evidence hashes:
python -B code/check_artifacts.py
\end{verbatim}

For a targeted neural replay, use
\path{report/support/interface_pinn/reevaluate.py}; it loads all six archived
checkpoints on CPU, solves the classical references and writes
\path{report/support/interface_pinn/comparison.json}. It does not retrain
networks. Then run \path{report/support/interface_pinn/make_figures.py}.
The training histories and GPU times are archived measurements, not outcomes
of the CPU replay. The single-network ablation has half the parameters and
is explicitly not parameter matched.

The strict bank-enabled validation passed all 11 tests, including 32
course-bank solves, as recorded in
\path{reviews/evidence/strict_tests_with_bank.txt} and
\path{reviews/evidence/course_bank_run.json}. To repeat it, set
\texttt{MTH321\_TEST\_BANK} to the instructor's local bank file and
\texttt{MTH321\_STRICT\_VALIDATION=1}, then run
\texttt{python -B -m unittest discover -s test -v}. The bank is external;
no skip is presented as a pass. The earlier independent measurements and
current targeted checks are distinguished in
\path{reviews/CHECKLIST_STATUS.md}.
